\documentclass[../preamble]{subfiles}
\begin{document}

\section{Natural Transformations}
\label{sec:natural-transformation}

% Typing macros
\providecommand{\caty}[1]{\mathsf{#1}}
\providecommand{\fun}[1]{\mathsf{#1}}
\providecommand{\nat}[1]{\mathsf{#1}}
\providecommand{\newmath}[1]{%
  {\setlength{\fboxsep}{0pt}%
   \mathchoice
     {\colorbox{defblue}{$\displaystyle #1$}}
     {\colorbox{defblue}{$\textstyle #1$}}
     {\colorbox{defblue}{$\scriptstyle #1$}}
     {\colorbox{defblue}{$\scriptscriptstyle #1$}}}}
\providecommand{\newterm}[1]{\textcolor{blue}{\emph{#1}}}

Let $\caty{C}$ and $\caty{D}$ be categories,
and let $\fun{F}, \fun{G} \colon \caty{C} \to \caty{D}$ be functors.
A \newterm{natural transformation}
$\newmath{\nat{\alpha} \colon \fun{F} \Rightarrow \fun{G}}$
consists of the following data:
\begin{enumerate}
  \item For each object $X \in \operatorname{Ob}(\caty{C})$,
    a morphism $\newmath{\nat{\alpha}_X} \colon \fun{F}(X) \to \fun{G}(X)$
    in $\caty{D}$,
    called the \newterm{component} of $\nat{\alpha}$ at $X$.
\end{enumerate}

subject to the \newterm{naturality condition}:
for every morphism $f \colon X \to Y$ in $\caty{C}$,
\begin{equation}
  \label{eq:naturality}
  \fun{G}(f) \circ \nat{\alpha}_X = \nat{\alpha}_Y \circ \fun{F}(f).
\end{equation}

This is expressed by the commutativity of the \newterm{naturality square}:
\begin{equation}
  \label{eq:naturality-square}
  \begin{tikzcd}
    \fun{F}(X) \arrow[r, "\nat{\alpha}_X"] \arrow[d, "\fun{F}(f)"']
      & \fun{G}(X) \arrow[d, "\fun{G}(f)"] \\
    \fun{F}(Y) \arrow[r, "\nat{\alpha}_Y"']
      & \fun{G}(Y)
  \end{tikzcd}
\end{equation}

\noindent
Reference: \textsc{Stacks Project} Tag \texttt{0015}.

\begin{definition}[Identity natural transformation]
  \label{def:id-nat-trans}
  The \newterm{identity natural transformation}
  $\newmath{\operatorname{id}_{\fun{F}}} \colon \fun{F} \Rightarrow \fun{F}$
  has components $(\operatorname{id}_{\fun{F}})_X = \operatorname{id}_{\fun{F}(X)}$
  for each $X \in \operatorname{Ob}(\caty{C})$.
  Naturality holds since both composites in~\eqref{eq:naturality-square}
  reduce to $\fun{F}(f)$.
\end{definition}

\begin{definition}[Vertical composition]
  \label{def:vertical-composition}
  Given natural transformations
  $\nat{\alpha} \colon \fun{F} \Rightarrow \fun{G}$
  and
  $\nat{\beta} \colon \fun{G} \Rightarrow \fun{H}$,
  their \newterm{vertical composite}
  $\newmath{\nat{\beta} \circ \nat{\alpha}} \colon \fun{F} \Rightarrow \fun{H}$
  has components
  \begin{equation}
    \label{eq:vertical-comp}
    (\nat{\beta} \circ \nat{\alpha})_X = \nat{\beta}_X \circ \nat{\alpha}_X
  \end{equation}
  for each $X \in \operatorname{Ob}(\caty{C})$.
\end{definition}

\begin{proposition}
  \label{prop:vertical-comp-nat}
  The vertical composite $\nat{\beta} \circ \nat{\alpha}$
  as defined in~\eqref{eq:vertical-comp} is a natural transformation
  $\fun{F} \Rightarrow \fun{H}$.
\end{proposition}

\begin{proof}
  For any morphism $f \colon X \to Y$ in $\caty{C}$,
  \[
    \fun{H}(f) \circ (\nat{\beta} \circ \nat{\alpha})_X
    = \fun{H}(f) \circ \nat{\beta}_X \circ \nat{\alpha}_X
    = \nat{\beta}_Y \circ \fun{G}(f) \circ \nat{\alpha}_X
    = \nat{\beta}_Y \circ \nat{\alpha}_Y \circ \fun{F}(f)
    = (\nat{\beta} \circ \nat{\alpha})_Y \circ \fun{F}(f),
  \]
  where the second and third equalities use naturality of $\nat{\beta}$
  and $\nat{\alpha}$ respectively.
\end{proof}

\begin{definition}[Horizontal composition / Godement product]
  \label{def:horizontal-composition}
  Given categories $\caty{C}$, $\caty{D}$, $\caty{E}$,
  functors $\fun{F}, \fun{F}' \colon \caty{C} \to \caty{D}$
  and $\fun{G}, \fun{G}' \colon \caty{D} \to \caty{E}$,
  and natural transformations
  $\nat{\alpha} \colon \fun{F} \Rightarrow \fun{F}'$
  and $\nat{\beta} \colon \fun{G} \Rightarrow \fun{G}'$,
  the \newterm{horizontal composite} (or \newterm{Godement product})
  $\newmath{\nat{\beta} \star \nat{\alpha}}
    \colon \fun{G} \circ \fun{F} \Rightarrow \fun{G}' \circ \fun{F}'$
  has components
  \begin{equation}
    \label{eq:horizontal-comp}
    (\nat{\beta} \star \nat{\alpha})_X
    = \nat{\beta}_{\fun{F}'(X)} \circ \fun{G}(\nat{\alpha}_X)
    = \fun{G}'(\nat{\alpha}_X) \circ \nat{\beta}_{\fun{F}(X)},
  \end{equation}
  for each $X \in \operatorname{Ob}(\caty{C})$.
  The two expressions in~\eqref{eq:horizontal-comp} are equal
  by naturality of $\nat{\beta}$ applied to $\nat{\alpha}_X$.
\end{definition}

\begin{remark}
  Vertical composition is associative and admits $\operatorname{id}_{\fun{F}}$
  as a two-sided unit.
  Together with the interchange law
  relating vertical and horizontal composition,
  these data make functors and natural transformations
  between two fixed categories into a category
  $\newmath{[\caty{C}, \caty{D}]}$
  (the \newterm{functor category}),
  and all small categories into a 2-category $\caty{Cat}$.
\end{remark}

\end{document}
